% page 216 du fac-simile (image p-215 du scan)
% bloc 160.0 x 242.0 mm ; marge gauche 8.0 mm
\pgc{-12.700mm}{0.000mm}{
\l{O8}
\l{}
\l{+}
\l{\sou{haraso.}\cel{3}I. Trupo ek kavaluli e kavalini di raso la maxim}
\l{\cel{3}bona, asemblita skope di genito. - II. Establisuro spe-}
\l{\cel{3}cala por la eduko e produkto di kavali. - F.}
\l{}
\l{\sou{haro}. (\sou{anat.}) Singla de la pili qui garnisas la kranio-pelo}
\l{\cel{3}(che la homo). - DE.}
\l{}
\l{\sou{harda}.\cel{2}I. (\sou{aludante kozo}). I. Qua tre rezistas lor tusheso.}
\l{\cel{3}II. Qua tre rezistas kontre entameso. - III. (\sou{aludant}e}
\l{\cel{3}\sou{persono}). Qua ne apertesas por la sentimenti di benigneso,}
\l{\cel{3}di humaneso, e c. - DE.}
\l{}
\l{\sou{haremo}. I. Che la Mohamedisti : apartmento di la homini.}
\l{\cel{3}II. Asemblitaro ek ta homini. - DEFIRS.}
\l{}
\l{\sou{haringo}. (\sou{zool.}) Mar-fisho de la familio \textquotedbl{}klupi\textquotedbl{} qua decen-}
\l{\cel{3}sas singla-yare, en trupi nekontebla, de la Oceano Glacia-}
\l{\cel{3}la Arktika. - L.\cel{1}\sou{Clupea harengus}. - DEFIS.}
\l{}
\l{\sou{harmoniar}. (\sou{netrans., kun)}\cel{1}. Havar sua parti tala, ke singla}
\l{\cel{3}de li konkordas kun la ceteri. (\sou{anke metaf.})\cel{2}DEFRS.}
\l{}
\l{\sou{harmoniko}. (\sou{matem.}) Proporciono ek tri nombri tala ke la ece-}
\l{\cel{3}so di la unesma ye la duesma relatas la eceso di la duesma}
\l{\cel{3}sur la triesma quale la duesma ye la unesma.}
\l{}
\l{\sou{harmoniko}. (\sou{muzik.}) I. Muzik-instrumento antiqua, de kalici}
\l{\cel{3}vitra, neegale plenigita ye aquo, e rotacanta, e qua, per la}
\l{\cel{3}froto di la bordo di la kalici da la fingri, bruisas vibre}
\l{\cel{3}- II. Muzik-instrumento en qua la kordi remplasesas da lami}
\l{\cel{3}vitra di qui la longeso ne esas inter-egala. - DEFIRS.}
\l{}
\l{\sou{harneso}. Equipo-kozi di sel-kavalo o jungo-kavalo.- DEFS.}
\l{}
\l{\sou{harpo}. Muzik-instrumento tri-angula, kun kordi vertikala, quin}
\l{\cel{3}onu pincas per la du manui e quin onu vibrigas. -DEFIRS.}
\l{}
\l{\sou{harpio}. I. (mitol.) Ento quan onu reprezentis kom havanta la}
\l{\cel{3}vizajo di muliero e la korpo di vulturo. - II. (\sou{metaf.})}
\l{\cel{3}Persono kapt-avida, muliero quaze rabiika. - DEFIRS.}
\l{}
\l{\sou{harpuno}. I. Instrumento fera,uzata por akrochar, pikar. - II.}
\l{\cel{3}Metal-squadro por ligar du peci ek konstrukturo. - DEFRS.}
\l{}
\l{\sou{hashisho}. Kanabo di India, di qua onu mastikas o fumas la fo-}
\l{\cel{3}lii sekigita. - DEFIRS.}
\l{}
\l{hastar. (netrans.,\cel{1}\sou{pri}). Ne disipar sua tempo (lor agar o}
\l{\cel{3}facar ulo). - DEF.}
\l{}
\l{hauo. Plug-instrumento fera kun ligno-mancho, por exkavar tero;}
\l{\cel{3}la fera parto esas larja, plata e tranchiva, kun aristo en}
\l{\cel{3}plano perpendikla ye la mancho. - DEF.}
}
